















\chapter{Computational Experiments} 
\label{computationalexp}
In this section, we describe our computational experiments with
the set partitioning formulation AUG-SPP-LRS and the branch-and-price 
algorithm described in Section \ref{SolnMethod}. We ran our experiments on a 
64-node Linux-based Beowulf cluster using condor. 


\section{Test Problems from \citet{Lin}}
To evaluate the performance of our algorithm, we solved the instances
presented in \citet{Lin}  which is solving single 
objective location routing and scheduling problem. 
\citet{Lin} present a real life data with 27 customers and 4 facilities 
and solve the instance using their heuristic algorithm. They also solve smaller 
instances that are subsets of the real data using branch and bound algorithm to show 
the difficulty of the LRSPs. Table \ref{CompLin} compares the results presented by \citet{Lin}
with the results of our branch and price algorithm. Since the instances
are not very difficult, to solve these instances we skipped the first step
and used only step 2 of our branch and price algorithm.  
As shown in table \ref{CompLin} part (a), branch and bound algorithm cannot solve the 
instances with 12 customers to the optimality. The computation time increases significantly
as the number of customer increases from 10 to 12. In the first 3 small instances,
our computation times are less than 1 minute and in 2 instances out of the first 3, 
our algorithm is faster than the heuristic presented in \citet{Lin}. For instance 4,
we can find the optimal cost which is less than the result given in  \citet{Lin} and 
the computation time is less than 8 minutes.   
\begin{table}[h!tbp]
  \begin{center}
    \caption{Comparing the results with \citet{Lin} and our BP algorithm} \label{CompLin}
    \begin {tabular}[h]{|c|c|c|c|c|}
      \hline 
{\bf Instances} & \multicolumn{2}{|c|}{{\bf Branch \& Bound}} & \multicolumn{2}{|c|}{{\bf Best Heurisitic}} \\\cline{2-5} 
\mbox{} & {\bf Obj.} & {\bf CPU (s)} & {\bf Obj.} & {\bf CPU (s)}  \\\hline
1: 4 depots, 10 custs.	&	309,817	&	1155	& 	309,817 & 	0.44 \\\hline
2: 4 depots, 10 custs.	&	309,808	&	982	&	309,808	&	0.49 	\\\hline
3: 4 depots, 12 custs.	&	312,036* & 	$>$ 10,000  & 	312,036 &       0.82 	\\\hline
4: 4 depots, 27 custs.	&       -       &	  -     &       625,752.5 &      6   	\\\hline
\multicolumn {5}{l}{* Best solution in 10000 CPU seconds in a Pentium III machine.} \\
\multicolumn{5}{c}{{\bf a. Lin et. al (2002) }} \\
\end{tabular}
  \begin {tabular}[h]{|c|c|c|}
      \hline 
{\bf Instances} & {\bf Obj.} &  {\bf CPU (s)} \\\hline
1: 4 depots, 10 custs.	& 309,816	&  0.93 \\\hline
2: 4 depots, 10 custs.	& 309,808    & 0.32	\\\hline
3: 4 depots, 12 custs.	& 312,036    & 0.13	\\\hline
4: 4 depots, 27 custs.	& 625,749	& 439.18 \\\hline
\multicolumn{3}{c}{{\bf b. 1 Step Branch and Price Algorithm }} 
\end{tabular}
%\end{footnotesize}
\end{center}
\end{table}
  
% \begin{table}
%\begin{sidewaystable}
%    \begin{center}
%    \begin{footnotesize}
%\caption{Comparing the results with \citet{Lin} and our algorithm}
%    \begin {tabular}[t]{|c|c|c|c|c|c|c|c|c|}
%      \hline 
%{\bf Instances} & \multicolumn{4}{|c|}{{\bf Lin et. al (1992) }} & \multicolumn{4}{|c|}{{\bf Our 2-Step Algorithm}} \\\cline{2-5}
%\mbox{} & \multicolumn{2}{|c|}{{\bf Branch \& Bound}} & \multicolumn{2}{|c|}{{\bf Best Heurisitic}} &  \multicolumn{4}{|c|}{\mbox{}} \\\cline{2-9} 
%\mbox{} & {\bf Obj.} & {\bf CPU (s)} & {\bf Obj.} & {\bf CPU (s)} & {\bf Obj.} & {\bf S1-CPU (s)} & {\bf S2-CPU (s)} & {\bf Total CPU} \\\hline
%1: 4 depots, 10 custs.	&	309,817	&	1155	& 	309,817 & 	0.44  & 309,816	&	0.89	&	0.82	& 1.71	\\\hline
%2: 4 depots, 10 custs.	&	309,808	&	982	&	309,808	&	0.49  & 309,807.83 &	0.38	&	0.18	& 0.56	\\\hline
%3: 4 depots, 12 custs.	&	312,036* & 	$>$ 10,000  & 	312,036 &       0.82  & 312,035.83 &	0.17	&       0.08	& 0.25	\\\hline
%4: 4 depots, 27 custs.	&       -       &	  -     &       625,752.5 &      6     &	625,749	&	2499.54	&	344.02	& 2843.56	\\\hline
%\multicolumn {9}{l}{* Best solution in 10000 CPU seconds in a Pentium III machine.} 
%\end{tabular}
%\end{footnotesize}
%\end{center}
%\end{table}
%\end{sidewaystable}

%%%%%%%%%%%%%%%%%%%%%%%%%%%%%%%%%%%%%%%%%%%%%%%%%%%%%%%%%%%%%%%%%%%%%%%%%%%%%%%%%%%%%%%%%%%%%%%%%%%%%%%%%%
%%%%%%%%%%%%%%%%%%%%%%%%%% COMMENT %%%%%%%%%%%%%%%%%%%%%%%%%%%%%%%%%%%%%%%%%%%%%%%%%%%%%%%%%%%%%%%%%%%%%%%%
%%%%%%%%%%%%%%%%%%%%%%%%%%%%%%%%%%%%%%%%%%%%%%%%%%%%%%%%%%%%%%%%%%%%%%%%%%%%%%%%%%%%%%%%%%%%%%%%%%%%%%%%%%
\begin{comment}
\begin{table}
%\begin{sidewaystable}
  \begin{center}
    %    \begin{footnotesize}
    \caption{Comparing the results with \citet{Lin} and our algorithm} \label{CompLin}
    \begin {tabular}[t]{|c|c|c|c|c|}
      \hline 
{\bf Instances} & \multicolumn{2}{|c|}{{\bf Branch \& Bound}} & \multicolumn{2}{|c|}{{\bf Best Heuristic}} \\\cline{2-5} 
\mbox{} & {\bf Obj.} & {\bf CPU (s)} & {\bf Obj.} & {\bf CPU (s)}  \\\hline
1: 4 depots, 10 custs.	&	309,817	&	1155	& 	309,817 & 	0.44 \\\hline
2: 4 depots, 10 custs.	&	309,808	&	982	&	309,808	&	0.49 	\\\hline
3: 4 depots, 12 custs.	&	312,036* & 	$>$ 10,000  & 	312,036 &       0.82 	\\\hline
4: 4 depots, 27 custs.	&       -       &	  -     &       625,752.5 &      6   	\\\hline
\multicolumn {5}{l}{* Best solution in 10000 CPU seconds in a Pentium III machine.} \\
\multicolumn{5}{c}{{\bf a. Lin et. al (2002) }} 
\end{tabular}
   \begin {tabular}[t]{|c|c|c|c|c|}
      \hline 
{\bf Instances} & {\bf Obj.} & {\bf S1-CPU (s)} & {\bf S2-CPU (s)} & {\bf Total CPU} \\\hline
1: 4 depots, 10 custs.	& 309,816	&	0.89	&	0.82	& 1.71	\\\hline
2: 4 depots, 10 custs.	& 309,807.83    &	0.38	&	0.18	& 0.56	\\\hline
3: 4 depots, 12 custs.	& 312,035.83    &	0.17	&       0.08	& 0.25	\\\hline
4: 4 depots, 27 custs.	& 625,749	&	2499.54	&	344.02	& 2843.56	\\\hline
\multicolumn{5}{c}{{\bf b. 2 Step Branch and Price Algorithm }} 
\end{tabular}
%\end{footnotesize}
\end{center}
\end{table}
%\end{sidewaystable}
\end{comment}

%\subsubsection{Test Problems Using Perl Instances}

\section{Test Problems Using Cordeau Instances}
\label{CordeauTest}
\subsection{Description of the Instances}
\label{ssec:DesCord}

To create an instance, we need to specify customer and candidate facility locations, 
customer demands, vehicle capacity, vehicle time limit, facility capacity, 
and facility and vehicle fixed and variable costs. 

We used four MDVRP benchmark problems (called $p01$, $p03$, $p07$, $p11$) 
developed by \citet{Cordeau2} to generate customer locations and customer demands. 
Each of the \citet{Cordeau2} instances that we selected has a different set of customer 
locations. We did not select instance $p02$ since the instance has the same customer 
locations with the instance $p01$. Similarly, the instances $p04$, $p05$, $p06$
and $p07$ and the instances $p08$, $p09$, $p10$ and $p11$ have the same customer locations.
We generated  8 set of instances with 25 and 40 customers. All of the instances have 
5 candidate facility locations. We included facility locations given for the related 
benchmark problem into the facility set and if the number of facilities is less than 
5 in the corresponding benchmark problem, we generated necessary number of facilities 
using a uniform distribution on the  
region where customers are located. We selected the facility capacity considering 
total demand and requiring at least two facilities to be in the solution.  
For each set of locations and demands, by using 2 different vehicle capacity 
values and 2 different vehicle time limits, we generated 32 (4x8) instances with 25 customers
and 32 instances  with 40 customers. The smallest vehicle capacity 
for an instance was chosen a value that covers at least one customer or two customers and 
a customer that has the largest demand. The smallest vehicle time limit was calculated 
as a value that is greater than two times the distance between every customer and the 
furthest facility.

Table \ref{probparam} presents parameter values used and explains
the generation of the instances.   
% the vehicle capacities, time limits and facility capacities 
%used for each group of instance.
\begin {table}[h]
\begin{center}
\caption {Generation of instances from \citet{Cordeau2} problems} \label{probparam}
\begin {tabular}[t]{llll} 
Inst. Gorup	&	Benchmark Prb.	&	LRSP Instance 	&	Customer Locations and Demands	\\\hline
1,2	&	p01	&	p01-f, p01-l	&	first and last 25 and 40 customers	\\
3,4,5	&	p03	&	p03-f, p03-m, p03-l	&	first, middle and last 25 and 40 customers	\\
6,7,8	&	p07	&	p07-f, p07-s, p07-t	&	first, second and third 25 and 40 customers	\\\hline
\multicolumn{4}{c}{{\bf a. Customer information }} 
\end{tabular}
\begin {tabular}[t]{lccccc} 
%\multicolumn {6}{c}{{\bf Table 1.} Problem Parameters}\\ \hline
Inst. Generated Using & Vcap1	& Vcap2	& TimeLimit1 &	TimeLimit2 & Fcap\\\hline
p01,p07	& 	50	&	70	&	140	&	160	&	250 \\
p03	&	60	&	80	&	140	&	160	&	250\\
p11	&	150	&	200	&	240	&	500	&	500\\\hline
\multicolumn{6}{c}{{\bf b. Parameter information }} 
\end{tabular}
\end{center}
\end {table} 

We named our instances based on the parameters. For example, the instance 
$p01-f25-v1t1$ was created using the {\em first} 25 customer locations of  
the \citet{Cordeau} instance called $p01$. The first possible vehicle capacity ($v1=50$) and time limit
($t1=140$) are used in the instance. Similarly, $p03-l25-v2t2$ was created using the {\em last} 
25 customer locations of the \citet{Cordeau} instance called $p03$. The second possible vehicle 
capacity ($v2=80$) and time limit ($t2=160$) are used in the instance.  

To estimate the facility fixed cost, we used the study done by \citet{Burlett}. Assuming
a 100,000 sq. ft. facility, we estimated that the daily fixed cost of a facility varies 
between \$1500 and \$2300 depending on the region in which the facility is located. We 
randomly generated 5 values between 1500 and 2300. We used the same facility fixed 
costs in all of our instances. We estimated that the 
daily fixed cost of a mid-size vehicle varies between \$200 and \$250 based on the
economic life analysis done by  Kenworth Truck Company \citeyearpar{Kenworth}. We used
the average value \$225 as the fixed cost of a vehicle in all of our instances.  
The operating cost of a vehicle is estimated to be \$60 per hour. To estimate the
operating cost, we considered the average vehicle speed, the average amount of 
fuel consumption, fuel cost and the ratio of our working hour limit to the real
life working hour limit, which is generally 8 hours. Table \ref{costparam} lists 
the cost parameters used.  
%To estimate
%the operating cost, we assumed that a vehicle travels 25 miles per hour, a vehicle 
%spends 1 gallon of fuel for every 6 miles, fuel costs vary between \$3.5 and \$6.0 
%per gallon and the ratio of our travel times to the real travel times is between
%140/480(8 hours) to 240/480. The last assumption finds the ratio of the working hour limit in the instance
%to 8 hours, which would be the real life working hour limit. We assumed that the fuel 
%cost for our travel times, which includes the  ratio of real time to our travel times, 
%is \$12. These assumptions results an operating cost of \$50 per hour. 


\begin {table}[h]
\begin{center}
\caption {Cost Parameters} \label{costparam}
\begin {tabular}[t]{ccccccc} \hline
%\multicolumn {6}{c}{{\bf Table 1.} Problem Parameters}\\ \hline
FC1 & FC2 & FC3 & FC4 & FC5 & VF & OC \\\hline
\$1615 /day  &	\$1845 /day  &	\$1781 /day  &	\$1975 /day  &	\$1753 /day  &   \$225 /day  & \$ 60 /hour  	 \\\hline
\end{tabular}
\end{center}
\end {table} 

The travel time, $t_{ik}$, between each pair of nodes is calculated using 
Euclidean distance calculations and then rounded to the nearest integer.  
We run the Floyd-Warshall Algorithm \citep{Ahuja} on the travel time
matrix to ensure that all travel times satisfy the triangle inequality.
In our test problems, we assume that the service time at each
node is zero. 
%We have used the update rule 
%that Lin et al used in \cite{Lin} that is assigning the half of the service time of 
%any node to the incoming arcs to that node and the other half to the
%outgoing arcs from that node. For example, let $s_{i}$ and $s_{k}$ 
%be the service times of nodes $i$ and $k$ respectively and $t_{ik}$ be 
%the travel time between nodes i and k. Basically, 
%$t'_{ik}=t_{ik}+\frac{1}{2}(s_{i}+s_{k})$ is the updated travel time. 
    
%\subsection{Results}
%In this section, we compare the performance of our heuristics to the optimal
%solution.  

\subsection{Computational Results}
In this section, we present the results of our computational experiments. 
We evaluate the strength of our set partitioning formulation, the performance of 
our pricing problem, and the performance of our 2-step branch-and-price algorithm. 

\subsubsection{Comparison of Alternative Formulations}
To compare the strength of the various formulations, we enumerated all possible 
pairings and solved the formulations using a standard branch-and-bound algorithm.
Since the number of feasible pairings is large, we are not able to enumerate
all columns of the instances with 25 customers and 30 customers. Instead, we constructed 
smaller instances from our test problems by taking the first 20 customer locations.
For the experiments in this part, we used the instances with 20 customers and
5 facilities. Since the constraints (\ref{spp_con1}) and (\ref{spp_con2}) are enough to 
determine the LRS problem, we construct our base formulation ($SPP_B$) from these
constraints. 

{\bf Experiment 1} To compare the strength of constraints (\ref{relation-Z-T}) and 
(\ref{relation}) introduced in section \ref{StrongForm} and to verify the 
value of their contribution to the lower bound of the LRS problem,
we construct two formulations: $SPP_1$ and $SPP_2$. $SPP_1$ is constructed
by adding the constraints (\ref{relation-Z-T}) to the $SPP_B$ and $SPP_2$ 
is constructed by adding the constraints (\ref{relation}) to the $SPP_B$.
Table \ref{formulations-gr1} shows the constraints (\ref{relation-Z-T}) and
(\ref{relation}), and lists the constraints in each formulation. 

\begin{table}[h]
\centering
\caption{LRS models for Experiment 1}\label{formulations-gr1}
\begin {tabular}[t]{|l|l|l|}  
\multicolumn{3}{l}{\mbox{}}\\
\multicolumn{3}{l}{Const.(\ref{relation-Z-T}): $Z_{jp} \leq T_{j} \ \ \forall j \in M, \forall p \in P_{j}$}\\
\multicolumn{3}{l}{\mbox{}}\\
\multicolumn{3}{l}{Const.(\ref{relation}): $\sum_{p \in P_{j}}a_{ipj}Z_{jp} \leq T_{j} \ \ \forall j \in M, \forall i \in N$}\\  
\multicolumn{3}{l}{\mbox{}}\\
\hline
$SPP_B$ & $SPP_1$ & $SPP_2$ \\\hline
1. SP Const.(\ref{spp_con1}) & 1. SP Const.(\ref{spp_con1}) & 1. SP Const.(\ref{spp_con1}) \\
2. Fac.Cap.Const.(\ref{spp_con2}) & 2. Fac.Cap.Const.(\ref{spp_con2}) & 2. Fac.Cap.Const.(\ref{spp_con2}) \\
\mbox{}  & 3. Relation Const.(\ref{relation-Z-T}) &  3. Relation Const.(\ref{relation})\\\hline
\end{tabular}
\end{table}

Table \ref{lpbounds-gr1} compares the LP relaxation bounds of $SPP_B$, $SPP_1$, $SPP_2$. In table 
\ref{lpbounds-gr1}, \%Improv columns for models $SPP_1$ and $SPP_2$ represent the percentage increase
in the LP relaxation bound with respect to the LP relaxation bound of the base model $SPP_B$. The results 
suggest that
the constraints (\ref{relation}) are better than the constraints (\ref{relation-Z-T}), since
the LP solution of the model $SPP_2$ is higher than the LP solution of the model $SPP_1$ in all
instances. The model $SPP_2$ improves the LP solution of the base model by 1.7\% to 17\% while 
the model $SPP_1$ only improves the LP solution of the base model by 0\% to 2\%.   

\begin{table}[h]
\centering
\caption{Results of Experiment 1}\label{lpbounds-gr1}
\begin {tabular}[t]{|c|c|c|c|c|c|c|c|}  
\hline
{\bf Data File} & {\bf \# of Col.} & {\bf IP} & {\bf $SPP_B$}& \multicolumn{2}{|c|}{{\bf $SPP_1$}} & \multicolumn{2}{|c|}{{\bf $SPP_2$}}\\ \cline {4-8}
\mbox{} &  \mbox{} &  \mbox{} &  {\bf LP} & {\bf LP} & {\bf \%Improv.} & {\bf LP} & {\bf \%Improv.} \\\hline
p01-f20	&	67558	&	4348	&	3273.822	&	3276.9	&	0.09	&	3471.73	&	6.045	\\\hline
p03-f20	&	52619	&	4394.33	&	3705.03	&	3705.07	&	0	&	3767.8	& 1.69		\\\hline
p03-l20	&	14779	&	4708.83	&	3552.272	&	3552.61	&	0.01	&	3755.31	&	5.716	\\\hline
p07-f20	&	68267	&	4357.17	&	2677.031	&	2687.63	&	0.4	&	2813.44	&	5.095	\\\hline
p11-f20	&	1950	&	5921.5	&	5291.045	&	5308.24	&	0.32	&	5849.5	&	10.555	\\\hline
p11-l20	&	898	&	8017.5	&	5865.493	&	5987.21	&	2.08	&	6855.28	&	16.875	\\\hline
\end{tabular}
\end{table} 

{\bf Experiment 2} After we observing that constraints (\ref{relation}) seem to be better than constraints 
(\ref{relation-Z-T}), we design experiment 2 to test the strength of constraint (\ref{lb-open-fac}) 
with respect to the base model and to test the combined affect of constraints (\ref{lb-open-fac}) and (\ref{relation}).
We construct two formulations: $SPP_3$ and $SPP_4$. $SPP_3$ is constructed
by adding the constraint (\ref{lb-open-fac}) to the base model $SPP_B$. 
$SPP_4$ is constructed by adding the constraints (\ref{relation}) to the $SPP_3$.
Table \ref{formulations-gr2} shows the constraint (\ref{lb-open-fac}) and lists the constraints in each formulation.   

\begin{table}[h]
\centering
\caption{LRS models for Experiment 2}\label{formulations-gr2}
\begin {tabular}[t]{|l|l|l|}  
\multicolumn{3}{l}{\mbox{}}\\
\multicolumn{3}{l}{Const.(\ref{lb-open-fac}): $\sum_{j \in M} T_{j} \geq nREQ$ }\\
\multicolumn{3}{l}{\mbox{}}\\
\hline
$SPP_B$ & $SPP_3$ & $SPP_4$ \\\hline
1. SP Cont.(\ref{spp_con1}) & 1. SP Cont.(\ref{spp_con1}) & 1. SP Cont.(\ref{spp_con1}) \\
2. Fac.Cap.Const.(\ref{spp_con2}) & 2. Fac.Cap.Const.(\ref{spp_con2}) & 2. Fac.Cap.Const.(\ref{spp_con2}) \\
\mbox{} & 3. Const.(\ref{lb-open-fac})  & 3. Const.(\ref{lb-open-fac}) \\
\mbox{} & \mbox{}  & 4. Relation Cont.(\ref{relation})\\\hline
\end{tabular}
\end{table}

Table \ref{lpbounds-gr2} presents the percentage improvement in the LP relaxations with
constraint (\ref{lb-open-fac}) and the percentage improvement in the LP relaxations with
constraints (\ref{lb-open-fac}) and (\ref{relation}). In table \ref{lpbounds-gr2}, the \%Improv. 
columns for models $SPP_2$, $SPP_3$ and $SPP_4$ represent the percentage increase
in the LP solutions with respect to the LP solution of the base model $SPP_B$. 
The results show that the LP bound improves by 6\% to 59\% with constraint 
(\ref{lb-open-fac}). When we compare models $SPP_2$ and $SPP_3$, we observe that
in most instances, constraint (\ref{lb-open-fac}) improves the LP relaxation bound 
more than the constraints (\ref{relation}) do. The two set of constraints, constraints (\ref{relation}) and
constraint (\ref{lb-open-fac}) together achieve larger improvement than 
constraint (\ref{lb-open-fac}) does alone in all instances. The LP relaxation
solution of the model $SPP_4$ is 10.5\% to 61\% better than the LP
solution of the base model. 

\begin{sidewaystable}
%\begin{table}[h]
\centering
\caption{Results of Experiment 2}\label{lpbounds-gr2}
\begin {tabular}[t]{|c|c|c|c|c|c|c|c|c|c|}  
\hline
{\bf Data File} & {\bf \# of Col.} & {\bf IP} & {\bf $SPP_B$}& \multicolumn{2}{|c|}{{\bf $SPP_2$}} & \multicolumn{2}{|c|}{{\bf $SPP_3$}} & \multicolumn{2}{|c|}{{\bf $SPP_4$}}\\ \cline {4-10}
\mbox{} &  \mbox{} &  \mbox{} &  {\bf LP} & {\bf LP} & {\bf \%Improv.} & {\bf LP} & {\bf \%Improv.} & {\bf LP} & {\bf \%Improv.} \\\hline
p01-f20	&	67558	&	4348	&	3273.822        & 3471.73	& 6.045  & 4251.74	&	29.87	&	4279.88	&	30.73	\\\hline
p03-f20	&	52619	&	4394.33	&	3705.03	        & 3767.8	& 1.69   &  4371.08	&	17.98	&	4383.24	&	18.31	\\\hline
p03-l20	&	14779	&	4708.83	&	3552.272	& 3755.31	& 5.716  & 4514.81	&	27.1	&	4656.04	&	31.07	\\\hline
p07-f20	&	68267	&	4357.17	&	2677.031	& 2813.44	& 5.095  & 4243.97	&	58.53	&	4312.98	&	61.11	\\\hline
p11-f20	&	1950	&	5921.5	&	5291.045	& 5849.5	& 10.555 & 5623.73	&	6.29	&	5849.5	&	10.55	\\\hline
p11-l20	&	898	&	8017.5	&	5865.493	& 6855.28	& 16.875 & 7198.17	&	22.72	&	7810.83	&	33.17	\\\hline
\end{tabular}
%\end{table} 
\end{sidewaystable}

Based on the experiments 1 and 2, we conclude that our valid inequalities improve the LP relaxation bound of the model, 
constraints (\ref{relation}) are better than  constraints (\ref{relation-Z-T}) and the $SPP_4$ formulation, which includes 
constraints (\ref{lb-open-fac}) and (\ref{relation}), is the best among other models since it has the highest 
lower bound in all instances. 


\qquad


\subsubsection{Comparison of ESPPRC Algorithm and Dynamic Programming}
\label{comp-espprc-dp}
**** VO cost = 50 is used in this experiments********************************* 

As we mentioned in Section \ref{Solve-pricing}, we can use a DP algorithm to  solve
a relaxed version of the pricing problem in which we allow customers to be visited
more than once. However, this weakens the LP relaxation bound. To evaluate empirically how much the bound
is weakened, we compared the LP relaxation bounds at the root node using the dynamic 
programming algorithm (DP) and the ESPPRC algorithm. Using the ESPPRC algorithm, we obtain
the true LP relaxation bound of the LRS model. With the DP algorithm, we obtain a lower 
bound on the LP relaxation bound since the solution potentially contains customers 
that are visited more than once. Table \ref{dp-espprc25} and   Table \ref{dp-espprc30} 
compare the LP relaxation bounds, total number of columns in the solution and the time 
spent to get the LP bounds. \%Gap represents the percentage gap between the LP relaxation 
solution obtained using the DP algorithm and the true LP relaxation solution (the LP relaxation obtained 
using ESPPRC algorithm). Basically, the \%Gap is calculated by dividing the difference between
the LP relaxation solutions by the true LP relaxation solution and multiplying by 100.  

The \%gap ranges from 0.98\% to 3.23\%. The average gap is 1.7\%  for the 
instances with 25 customers and 1.81\% for the instances with 30 customers. The 
ESPPRC algorithm is preferable to the DP algorithm in all of the $p11$ instances. 
In general (except $p11$ instances), ESPPRC algorithm takes more time than the DP algorithm,
which is the price we have to pay for a better lower bound. Although in one instance 
(p01-f25-v2t2) the ESPPRC algorithm took a long time, generally the computation times are 
not too long for 25 customer instances. Therefore, the ESPPRC algorithm is preferable
to the DP algorithm. However, in 9 instances out of 24 with 30 customers, the ESPPRC 
algorithm took more than 2 hours. Hence, with larger instances, depending on the difficulty of the 
instance,  we may search for ways to improve our current ESPPRC algorithm.
%using the DP algorithm in combination with the ESPPRC algorithm 
%to solve the root node sooner.  
  
\begin{sidewaystable}
%\begin{table}
\centering
\caption {Root node solutions using the DP algorithm and the ESPPRC algorithm (25 customer instances)} \label{dp-espprc25}
\begin {tabular}[b]{|c|c|c|c|c|c|c|c|}
%\multicolumn {}{c}{{\bf Table 4.} Root node solutions using ESPPRC and DP}\\\hline
\hline
{\bf Data File} & \multicolumn{3}{|c|}{{\bf DP Algorithm Root Node}} &  \multicolumn {3}{|c|}{{\bf ESPPRC Algorithm Root node}} & {\bf $\%$Gap}\\ \cline {2-7}
\mbox{} & {\bf LP Relax.} & {\bf $\#$ of Columns} & {\bf	Time (sec.)} &	{\bf LP Relax.}	& {\bf $\#$ of Columns}	& {\bf Time (sec.)} & \mbox{} \\\hline
p01-f25-v1t1	&	4352.05	&	7557	&	131.74	&	4404.3	&	8889	&	611.5	&	1.19	\\\hline
p01-f25-v1t2	&	4268.98	&	7667	&	146.49	&	4321.41	&	9317	&	3536.46	&	1.21	\\\hline
p01-f25-v2t1	&	4218.63	&	11956	&	203.07	&	4266.38	&	13292	&	2278.57	&	1.12	\\\hline
p01-f25-v2t2	&	4148.12	&	12192	&	274.51	&	4192.23	&	13697	&	18434.96	&	1.05	\\\hline \hline
p03-f25-v1t1	&	4651.47	&	7814	&	150.67	&	4699.51	&	8705	&	244.01	&	1.02	\\\hline
p03-f25-v1t2	&	4540.84	&	8297	&	221.43	&	4582.2	&	9347	&	1027.37	&	0.9	\\\hline
p03-f25-v2t1	&	4446.23	&	11921	&	237.21	&	4492.06	&	12731	&	619.33	&	1.02	\\\hline
p03-f25-v2t2	&	4349.95	&	12762	&	269.98	&	4404.59	&	13850	&	3577.31	&	1.24	\\\hline \hline
p03-l25-v1t1	&	4605.93	&	8589	&	142.26	&	4682.91	&	9093	&	118.34	&	1.64	\\\hline
p03-l25-v1t2	&	4471.96	&	10151	&	148.15	&	4562.92	&	11061	&	750.31	&	1.99	\\\hline
p03-l25-v2t1	&	4428.19	&	9808	&	216.17	&	4526.2	&	10440	&	334.96	&	2.17	\\\hline
p03-l25-v2t2	&	4320.83	&	11791	&	295.63	&	4416.08	&	12855	&	1538.03	&	2.16	\\\hline \hline
p07-f25-v1t1	&	4499.74	&	10275	&	121.52	&	4560.82	&	10979	&	210.08	&	1.34	\\\hline
p07-f25-v1t2	&	4401.35	&	11132	&	147.76	&	4457.51	&	12049	&	1135.77	&	1.26	\\\hline
p07-f25-v2t1	&	4310.75	&	12133	&	187.29	&	4408.23	&	12898	&	526.07	&	2.21	\\\hline
p07-f25-v2t2	&	4232.34	&	13260	&	222.19	&	4303.73	&	14368	&	4139.74	&	1.66	\\\hline \hline
p11-f25-v1t1	&	7515.63	&	1758	&	205.15	&	7601.51	&	1883	&	0.94	&	1.13	\\\hline
p11-f25-v1t2	&	7444	&	2335	&	258.47	&	7521.76	&	2540	&	1.59	&	1.03	\\\hline
p11-f25-v2t1	&	7116.49	&	2495	&	362.47	&	7371.89	&	2629	&	1.53	&	3.46	\\\hline
p11-f25-v2t2	&	7048.41	&	3396	&	408.73	&	7254.45	&	3606	&	3.12	&	2.84	\\\hline \hline
p11-l25-v1t1	&	8047.81	&	1086	&	192.4	&	8178.02	&	1153	&	0.52	&	1.59	\\\hline
p11-l25-v1t2	&	7876.72	&	1497	&	213.78	&	8036.66	&	1563	&	0.84	&	1.99	\\\hline
p11-l25-v2t1	&	7447.09	&	1490	&	247.02	&	7633.51	&	1565	&	0.81	&	2.44	\\\hline
p11-l25-v2t2	&	7341.22	&	2087	&	285.32	&	7579.66	&	2168	&	1.34	&	3.15	\\\hline \hline
\end{tabular}
\end{sidewaystable} 
%\end{table}

  
\begin{sidewaystable}
%\begin{table}
\centering
\caption {Root node solutions using the DP algorithm and the ESPPRC algorithm (30 customer instances)} \label{dp-espprc30}
\begin {tabular}[b]{|c|c|c|c|c|c|c|c|}
%\multicolumn {}{c}{{\bf Table 4.} Root node solutions using ESPPRC and DP}\\\hline
\hline
{\bf Data File} & \multicolumn{3}{|c|}{{\bf DP Algorithm Root Node}} &  \multicolumn {3}{|c|}{{\bf ESPPRC Algorithm Root node}} & {\bf $\%$Gap}\\ \cline {2-7}
\mbox{} & {\bf LP Relax.} & {\bf $\#$ of Columns} & {\bf	Time (sec.)} &	{\bf LP Relax.}	& {\bf $\#$ of Columns}	& {\bf Time (sec.)} & \mbox{} \\\hline

p01-f30-v1t1	&	4773.67	&	13840	&	222.04	&	4821.07	&	15203	&	2901.33	&	0.98	\\\hline
p01-f30-v1t2	&	4666.71	&	14267	&	253.38	&	4715.3	&	15859	&	32779.86	&	1.03	\\\hline
p01-f30-v2t1	&	4514.93	&	20846	&	378.32	&	4582.27	&	22312	&	11770.67	&	1.47	\\\hline
p01-f30-v2t2	&	4414.87	&	21405	&	450.19	&	4482.2	&	23361	&	211885.55	&	1.5	\\\hline \hline
p03-f30-v1t1	&	6403.69	&	13332	&	254.71	&	6467.35	&	14459	&	2194.21	&	0.98	\\\hline
p03-f30-v1t2	&	6298.48	&	13851	&	233.95	&	6370.12	&	15165	&	28130.43	&	1.12	\\\hline
p03-f30-v2t1	&	6183.48	&	20904	&	404.95	&	6263.17	&	22392	&	11978.05	&	1.27	\\\hline
p03-f30-v2t2	&	6099.61	&	21912	&	426.91	&	6193.52	&	23418	&	197927.28	&	1.52	\\\hline \hline
p03-l30-v1t1	&	4786.72	&	14751	&	297.23	&	4896.04	&	15678	&	788.66	&	2.23	\\\hline
p03-l30-v1t2	&	4652.91	&	17388	&	354.24	&	4756.95	&	18521	&	5342.55	&	2.19	\\\hline
p03-l30-v2t1	&	4555.76	&	17407	&	456.01	&	4679.23	&	18641	&	2402.75	&	2.64	\\\hline
p03-l30-v2t2	&	4453.71	&	20749	&	541.96	&	4562.29	&	22270	&	20511.51	&	2.38	\\\hline \hline
p07-f30-v1t1	&	4706.19	&	17515	&	218.28	&	4765.63	&	18968	&	1863.84	&	1.25	\\\hline
p07-f30-v1t2	&	4597.22	&	18563	&	247.78	&	4654.19	&	20190	&	10036.58	&	1.22	\\\hline
p07-f30-v2t1	&	4474.19	&	21332	&	345.67	&	4558.21	&	22656	&	4517.61	&	1.84	\\\hline
p07-f30-v2t2	&	4380.11	&	22783	&	377.95	&	4453.2	&	24522	&	84685.95	&	1.64	\\\hline \hline
p11-f30-v1t1	&	7774.24	&	2055	&	356.58	&	7916.08	&	2169	&	1.8	&	1.79	\\\hline
p11-f30-v1t2	&	7681.77	&	2809	&	541.54	&	7791.1	&	2985	&	3.06	&	1.4	\\\hline
p11-f30-v2t1	&	7307.09	&	2628	&	526.21	&	7531.81	&	2834	&	4	&	2.98	\\\hline
p11-f30-v2t2	&	7239.85	&	3670	&	738.31	&	7469.84	&	3978	&	8.29	&	3.08	\\\hline \hline
p11-l30-v1t1	&	9804.77	&	2062	&	350.01	&	9929.75	&	2157	&	1.33	&	1.26	\\\hline
p11-l30-v1t2	&	9657.35	&	2753	&	388.93	&	9837.09	&	2888	&	2.9	&	1.83	\\\hline
p11-l30-v2t1	&	9386.09	&	2727	&	440.88	&	9645.66	&	2850	&	2.5	&	2.69	\\\hline
p11-l30-v2t2	&	9112.56	&	3727	&	476.42	&	9416.78	&	3961	&	6.51	&	3.23	\\\hline \hline

\end{tabular}
\end{sidewaystable} 
%\end{table}

   
\subsubsection {Results of 2-Step Branch-and-Price Algorithm}

This experiment is designed to evaluate the performance of the 2-Step branch and price algorithm.
In this experiment, the first step of the algorithm
is run with 3 CPU hour time limit and the second step is run with 6 CPU hour time limit. Computational results
report LP solutions, integrality gaps, IP solutions which are the optimal solutions if the integrality gap is
0, CPU times of each step and some solution information that includes  number of open facilities and number 
of vehicles used. 
 
{\bf Instances with 25 Customers}: Table \ref{25Cust_res} presents the computational results for 32 instances 
(explained in table \ref{probparam}) with 25 customers and 5 facilities.  As seen from the table  
 \ref{25Cust_res} we can find optimal solutions of 31 instances out of 32 instances in less than 4 hours. 
For one instance, the algorithm terminates due to the time limit and results an integer 
solution with an integrality gap of 0.11\%. 

%-Number of nodes vs time!
%{\bf Comparing Effect of vehicle capacity timelimit}
%++++++++++++++++++++++++++++++++++++++\\
%PUT a graph showing the time wrt vehicle capacity and timelimit \\
%++++++++++++++++++++++++++++++++++++++++

{\bf Instances with 40 Customers}:  
We also evaluate the performance of 2-Step BP Algorithm for instances (explained in table \ref{probparam}) with 40
customers and 5 facilities. Problem becomes more difficult when the number of customers increases. 
Table \ref{40Cust_res2} presents the computational results for the instances with 40 customers and 5 facilities. 
In 28 instances out of 32, the algorithm finds the LP solution and an IP solution for the LRSP. In 4 instances 
(p01-f40-v2t2, p03-f40-v2t2, p03-l40-v2t2, p07-t40-v2t2),
the algorithm could not find the exact LP solution within the time limit. However, in 2 of these 4 instances,
we could find a lower bound for the LP solution using the solution of exact pricing problem using the
formula \ref{LPlb} (see \citet{Lubbecke}):
\begin{equation}
z_{LB}(LP) = z + \sum_{j \in M}\hat{c}_j \label{LPlb}
\end{equation}
where $z$ is the current LP solution of the RMP and $\hat{c}_j$ is the reduced cost of the 
column for facility $j$ obtained from the exact pricing problem of facility $j$.

The instances (called px-v2t2) with larger time limit and larger vehicle capacity are
very difficult, to improve the integer solution obtained from the algorithm, we used a simple fact: Every feasible 
solution for instance  px-v1t1 is also feasible for px-v1t2, px-v2t1 and px-v2t2, and  every feasible 
solution for instances  px-v2t1 and px-v2t1 is feasible for px-v2t2. Note that in px-v1t1, px-v1t2,
px-v2t1 and px-v2t2 every parameter is same except the time limit and vehicle capacity. In px-v1t2 and
px-v2t2 the time limit is larger than the time limit used in px-v1t1 and in px-v2t1.In px-v2t1 and
px-v2t2 the vehicle capacity is larger than the vehicle capacity used in px-v1t1 and in px-v1t2.
At the end of algorithm, the condition is checked for px-v2t2 instances and if it applies, the upper bounds 
obtained from the algorithm is updated. In couple instances in group px-v2t2, because of the difficulty of the
pricing problems, the upper bound obtained 
at the end of the algorithm is worse than the best upper bound obtained from instances px-v1t1, px-v1t2
and px-v2t1. 

As shown in graph \ref{Gapnumber}, in 14 instances the algorithm can find optimal
solutions or   upper bounds with integrality gaps less than 3\%, in 7 instances gaps are  
between 3\% and 5\%, in 6 instances the gaps are between 5\% and 7\% and
in 3 instances the gaps are between 7\% and 10.5\%. 

\begin{figure}[h]
  \centering
  \includegraphics[scale=0.7]{figs/gapnumber.eps}
  %\fbox{ \includegraphics[scale=0.5]{figure/pair2.eps}}
  \caption{Number of instances with 40 customers with gaps in each region} \label{Gapnumber}
\end{figure}


\newpage
\begin{table}[h!tbp]
\centering
\caption{Results for instances with 25 customer 5 facilities}\label{25Cust_res}
\begin {tabular}[t]{|c|c|c|c|c|c|c|c|c|}  
\hline
{\bf Data File} & {\bf LP} & {\bf IP$^1$} & {\bf Gap} & {\bf Nodes} & \multicolumn{3}{|c|}{\bf CPU Time (hr)} & {\bf Soln: \# of} \\\cline{6-8}  
\mbox{} & \mbox{} & \mbox{} & \mbox{} & \mbox{} &  {\bf Step 1 } & {\bf Step 2} & {\bf Total} & {\bf Fac. - Veh. } \\\hline
p01-f25-v1t1	&	4504.4	&	4699	&	0	&	27	&	1.22	&	0.13	&	1.36	&	2	-	4	\\
p01-f25-v1t2	&	4425.79	&	4511	&	0	&	103	&	1.71	&	0.75	&	2.46	&	2	-	4	\\
p01-f25-v2t1	&	4366.12	&	4425	&	0	&	45	&	1.5	&	0.42	&	1.93	&	2	-	3	\\
p01-f25-v2t2	&	4295.4	&	4425	&	0	&	39	&	1.03	&	1.83	&	2.87	&	2	-	3	\\
p01-l25-v1t1	&	4849.78	&	4899	&	0	&	21	&	0.14	&	0.04	&	0.19	&	2	-	5	\\
p01-l25-v1t2	&	4725.39	&	4815	&	0	&	81	&	0.42	&	0.37	&	0.79	&	2	-	4	\\
p01-l25-v2t1	&	4668.98	&	4759	&	0	&	147	&	0.54	&	0.53	&	1.07	&	2	-	4	\\
p01-l25-v2t2	&	4537.07	&	4602	&	0	&	37	&	1.44	&	0.99	&	2.43	&	2	-	5	\\
p03-f25-v1t1	&	4809.96	&	5062	&	0	&	299	&	1.43	&	0.41	&	1.84	&	2	-	5	\\
p03-f25-v1t2	&	4701.14	&	4837	&	0	&	105	&	1.28	&	0.22	&	1.5	&	2	-	4	\\
p03-f25-v2t1	&	4606.09	&	4758	&	0	&	283	&	1.36	&	0.35	&	1.7	&	2	-	4	\\
p03-f25-v2t2	&	4502.54	&	4742	&	0	&	45	&	1.55	&	1	&	2.54	&	2	-	4	\\
p03-m25-v1t1	&	4697.37	&	4785	&	0	&	263	&	1.26	&	0.23	&	1.49	&	2	-	4	\\
p03-m25-v1t2	& 4578.24	&	4785	&	0.11	&	2401	&	3	&	6	&	9	&	2	-	4	\\
p03-m25-v2t1	&	4485.79	&	4689	&	0	&	59	&	1.33	&	0.44	&	1.77	&	2	-	4	\\
p03-m25-v2t2	&	4396.12	&	4479	&	0	&	9	&	1.4	&	1.72	&	3.12	&	2	-	3	\\
p03-l25-v1t1	&	4795.3	&	4842	&	0	&	15	&	0.12	&	0.04	&	0.16	&	2	-	4	\\
p03-l25-v1t2	&	4670.24	&	4803	&	0	&	37	&	0.29	&	0.24	&	0.53	&	2	-	4	\\
p03-l25-v2t1	&	4620.67	&	4741	&	0	&	43	&	0.25	&	0.19	&	0.43	&	2	-	4	\\
p03-l25-v2t2	&	4507.62	&	4655	&	0	&	39	&	0.67	&	0.88	&	1.55	&	2	-	3	\\
p07-f25-v1t1	&	4642.59	&	4761	&	0	&	51	&	0.21	&	0.16	&	0.37	&	2	-	4	\\
p07-f25-v1t2	&	4544.45	&	4685	&	0	&	69	&	1.33	&	0.87	&	2.2	&	2	-	3	\\
p07-f25-v2t1	&	4485.56	&	4703	&	0	&	161	&	1.53	&	0.63	&	2.16	&	2	-	4	\\
p07-f25-v2t2	&	4388.33	&	4482	&	0	&	125	&	1.03	&	1.04	&	2.08	&	2	-	3	\\
p07-s25-v1t1	&	5000.13	&	5098	&	0	&	257	&	0.07	&	0.05	&	0.11	&	2	-	5	\\
p07-s25-v1t2	&	4852.2	&	4889	&	0	&	75	&	0.26	&	0.11	&	0.37	&	2	-	4	\\
p07-s25-v2t1	&	4739.42	&	4787	&	0	&	23	&	0.08	&	0.04	&	0.12	&	2	-	4	\\
p07-s25-v2t2	&	4609.82	&	4781	&	0	&	85	&	0.43	&	0.36	&	0.8	&	2	-	4	\\
p07-t25-v1t1	&	4836.55	&	4898	&	0	&	29	&	0.12	&	0.1	&	0.22	&	2	-	4	\\
p07-t25-v1t2	&	4703.06	&	4886	&	0	&	467	&	1.92	&	1.35	&	3.27	&	2	-	4	\\
p07-t25-v2t1	&	4646.18	&	4767	&	0	&	247	&	0.42	&	0.59	&	1.02	&	2	-	4	\\
p07-t25-v2t2	&	4513.08	&	4592	&	0	&	121	&	1.45	&	2.03	&	3.47	&	2	-	3	\\\hline
\multicolumn{9}{l}{$^1$Optimal or best IP found by the algorithm within the time limit.}\\
\end{tabular}
\end{table} 


\newpage
\begin{table}[h!tbp]
\centering
\caption{{\bf Updated} Results for instances with 40 customer 5 facilities}\label{40Cust_res2}
\begin {tabular}[t]{|c|c|c|c|c|c|c|c|c|}  
\hline
{\bf Data File} & {\bf LP} & {\bf IP$^1$} & {\bf Gap} & {\bf Nodes} & \multicolumn{3}{|c|}{\bf CPU Time (hr)} & {\bf Soln: \# of} \\\cline{6-8}  
\mbox{} & \mbox{} & \mbox{} & \mbox{} & \mbox{} &  {\bf Step 1 } & {\bf Step 2} & {\bf Total} & {\bf Fac. - Veh. } \\\hline
p01-f40-v1t1	&	6976.02	&	7170	&	0.65\%	&	78	&	3	&	6	&	9	&	5	-	3	\\
p01-f40-v1t2	&	6826.16	&	7170	&	4.28\%	&	8	&	3	&	6.03	&	9.04	&	5	-	3	\\
p01-f40-v2t1	&	6700.21	&	7064	&	4.58\%	&	7	&	3	&	6.11	&	9.11	&	5	-	3	\\
p01-f40-v2t2	&	{\bf 6535.65$^{2}$}	&	6984	&	{\bf 6.41\%}	&	1	&	3	&	6.17	&	9.18	&	5	-	3	\\
p01-l40-v1t1	&	6899.63	&	7159	&	2.92\%	&	315	&	3	&	6	&	9	&	6	-	3	\\
p01-l40-v1t2	&	6757.39	&	6934	&	1.25\%	&	15	&	2.98	&	6	&	8.99	&	5	-	3	\\
p01-l40-v2t1	&	6616.66	&	6830	&	2.06\%	&	11	&	3	&	6.12	&	9.13	&	5	-	3	\\
p01-l40-v2t2	&	6513.61	&	{\bf 6830$^{3}$}	&	{\bf 4.36\%}	&	2	&	3	&	7.17	&	10.17	&	6	-	3	\\
p03-f40-v1t1	&	8697.41	&	8780	&	0	&	67	&	2.68	&	4.12	&	6.79	&	5	-	4	\\
p03-f40-v1t2	&	8578.62	&	8753	&	1.78\%	&	7	&	2.95	&	6.35	&	9.3	&	5	-	4	\\
p03-f40-v2t1	&	8439	&	8440	&	0.01\%	&	6	&	0.33	&	6.59	&	6.92	&	4	-	3	\\
p03-f40-v2t2	&	{\bf No Soln}	&	8438	&	\mbox{}	&	1	&	2.01	&	6.15	&	8.15	&	4	-	4	\\
p03-m40-v1t1	&	7050.88	&	7722	&	6.16\%	&	723	&	3	&	6	&	9	&	8	-	3	\\
p03-m40-v1t2	&	6907.66	&	7544	&	7.14\%	&	17	&	3	&	6.07	&	9.07	&	6	-	3	\\
p03-m40-v2t1	&	6793.88	&	7479	&	7.77\%	&	19	&	3	&	6.05	&	9.05	&	5	-	3	\\
p03-m40-v2t2	&	6655.5	&	7430	&	10.42\%	&	2	&	3	&	6.76	&	9.76	&	6	-	3	\\
p03-l40-v1t1	&	6920.38	&	7386	&	4.81\%	&	55	&	3	&	6.02	&	9.02	&	6	-	3	\\
p03-l40-v1t2	&	6771.56	&	7212	&	5.62\%	&	3	&	3	&	6.1	&	9.1	&	8	-	4	\\
p03-l40-v2t1	&	6662.92	&	7121	&	5.68\%	&	10	&	3	&	6	&	9.01	&	6	-	3	\\
p03-l40-v2t2	&	{\bf 6533.35$^{2}$}	&	7018	&	{\bf 6.91\%}	&	1	&	3	&	6	&	9	&	4	-	2	\\
p07-f40-v1t1	&	7053.07	&	7167	&	0	&	105	&	3	&	4.87	&	7.87	&	5	-	3	\\
p07-f40-v1t2	&	6906.08	&	7340	&	5.24\%	&	19	&	3	&	6.04	&	9.04	&	6	-	3	\\
p07-f40-v2t1	&	6745.5	&	6977	&	2.54\%	&	22	&	3	&	6.04	&	9.04	&	5	-	3	\\
p07-f40-v2t2	&	6629.41	&	{\bf 6977$^{3}$}	&	{\bf 4.98\%}	&	2	&	3	&	6.9	&	9.9	&	5	-	3	\\
p07-s40-v1t1	&	7256.31	&	7343	&	0	&	73	&	1.88	&	0.75	&	2.63	&	6	-	3	\\
p07-s40-v1t2	&	7091.11	&	7117	&	0	&	23	&	1.19	&	1.7	&	2.89	&	5	-	3	\\
p07-s40-v2t1	&	6924.49	&	7209	&	2.54\%	&	212	&	3	&	6	&	9	&	6	-	3	\\
p07-s40-v2t2	&	6762.76	&	7022	&	2.85\%	&	17	&	3	&	6	&	9	&	5	-	3	\\
p07-t40-v1t1	&	6940.52	&	7255	&	3.45\%	&	43	&	3	&	6.06	&	9.06	&	6	-	3	\\
p07-t40-v1t2	&	6806.52	&	7058	&	3.56\%	&	2	&	3	&	6.19	&	9.19	&	5	-	3	\\
p07-t40-v2t1	&	6649.81	&	6663	&	0.19\%	&	4	&	3	&	6.12	&	9.12	&	4	-	3	\\
p07-t40-v2t2	&	{\bf No Soln}	&	{\bf 6663$^{3}$}	&	\mbox{}	&	1	&	3	&	6	&	9	&	5	-	3	\\\hline
\multicolumn{9}{l}{$^{1}$ Optimal or best IP found by the algorithm within the time limit.}\\
\multicolumn{9}{l}{$^{2}$ Lower bound for the LP solution.}\\
\multicolumn{9}{l}{$^{3}$ IP solution is updated using the instances with smaller vehicle capacity and time limit.}\\

\end{tabular}
\end{table}  

{\bf Effect of Facility Fixed Costs:} With a subset of  instance group with 40 customers, we measured the
effect of facility fixed costs on the difficulty of the problem. We used 7 instances in this experiment. 
We investigated 4 different cases: 1. Base case which is the results in table \ref{40Cust_res2} in
which the facility fixed costs are generated using a uniform distribution between 1500 and 2300 as explained
in section \ref{ssec:DesCord}, 2. Half of base case in which the half of the facility fixed costs used in base case
are used, 3. Same fixed costs case in which the average fixed cost \$1900 is used for every facility, 
4. 0 Fixed cost case in which the facility fixed costs are zero. Table \ref{CompFacFixedC} shows
the results of the experiment. 
 
\begin{table}[h!tbp]
\centering
\caption{Effect of Facility Fixed Cost}\label{CompFacFixedC}
\begin {tabular}[t]{|c|c|c|c|c|c|c|c|c|}  
\hline 
Instance & \multicolumn{2}{|c|}{1. Base Case} &  \multicolumn{2}{|c|}{2. Half of Base Case}	&  \multicolumn{2}{|c|}{3. Same Fixed Costs Case}	&  \multicolumn{2}{|c|}{4. Zero Fixed Cost} \\\cline{2-9}			
\mbox{} & Gap &	TCPU (hr)	& Gap &	TCPU (hr)	& Gap & TCPU (hr)	& Gap & TCPU (hr)\\\hline
p03-f40-v1t1	&	0.00\%	&	6.79	&	0.00\%	&	8.15	&       0.00\%	&	7.41	&	0.00\%	&	4.84	\\
p03-f40-v1t2	&	1.78\%	&	9.3	&	1.95\%	&	9.16	&	0.00\%	&	5.96	&	0.00\%	&	4.81	\\
p03-f40-v2t1	&	0.01\%	&	6.92	&	4.45\%	&	9.36	&	0.08\%	&	7.53	&	0.44\%	&	7.58	\\
p07-f40-v1t1	&	0.00\%	&	7.87	&	0.00\%	&	5.25	&	0.00\%	&	7.41	&	0.00\%	&	6.16	\\
p07-f40-v1t2	&	5.24\%	&	9.04	&	4.92\%	&	9.04	&	3.19\%	&	9.04	&	12.69\%	&	9.17	\\
p07-f40-v2t1	&	2.54\%	&	9.04	&	4.18\%	&	9.02	&	2.68\%	&	9.06	&	0.00\%	&	6.07	\\
p07-f40-v2t2	&	4.98\%	&	9.9	&	7.18\%	&	9.78	&	4.65\%	&	9.32	&	5.95\%	&	10.11	\\\hline
\end{tabular}
\end{table}  

**** There is not a significant pattern! What should be our comment? *****

{\bf Comparing Vehicle Fixed Costs:} To see the effect of the magnitude of the vehicle fixed cost,
we designed an experiment that compares 3 cases using 8 difficult instances with 40 customers and using 
parameters Vcap2 and timelimit2 (px-v2t2). Note that the instance px-v2t2 is the most difficult instance
in group px (px-v1t1, px-v1t2, px-v2t1, px-v2t2). The base case is formed with the results in table
\ref{40Cust_res2} in which we used \$225 as a vehicle cost. We solved the instance group
for this experiment with vehicle fixed cost of \$110 (Case 2) and with 0 vehicle fixed cost
(Case 3).  In this experiment, since we only solved a subset of our instance group, 
we did not update the upper bounds for px-v2t2 instances using the IP solutions
of instances px-v1t1, px-v1t2 or px-v2t1.  
 Table \ref{CompVehFixedC} shows the results of the experiment. The results support that
increasing the vehicle fixed cost increases the difficulty of the problem. While in the base case,
the algorithm can not find the exact LP solution for 4 instances, in case 2 and 3, LP solutions can be 
found for all instances. In addition, the gaps for all instances in case 2 is less than the gaps in
base case. And the gaps for all instances in case 3 is less than the gaps in
case 3. In case 3 with zero fixed cost, the algorithm can almost find the optimal integer
solutions.  


\begin{table}[h!tbp]
\centering
\caption{Effect of Vehicle Fixed Cost}\label{CompVehFixedC}
\begin {tabular}[t]{|c|c|c|c|c|c|c|}  
\hline 
Instance & \multicolumn{2}{|c|}{Base Case: VF=\$225} &  \multicolumn{2}{|c|}{Case 2: VF=\$110}	&  \multicolumn{2}{|c|}{Case 3: VF=\$0} \\\cline{2-7}
\mbox{} & Exact LP found? & Gap & Exact LP found? & Gap & Exact LP found? & Gap \\\hline

p01-f40-v2t2	&	No	&	-	&	yes	&	4.07\%	&	yes	&	0.53\%	\\
p01-l40-v2t2	&	yes	&	9.10\%	&	yes	&	0.67\%	&	yes	&	0.03\%	\\
p03-f40-v2t2	&	No	&	-	&	yes	&	0.14\%	&	yes	&	0.02\%	\\
p03-m40-v2t2	&	yes	&	10.42\%	&	yes	&	6.98\%	&	yes	&	0.23\%	\\
p03-l40-v2t2	&	No	&	-	&	yes	&	4.79\%	&	yes	&	0.04\%	\\
p07-f40-v2t2	&	yes	&	5.29\%	&	yes	&	1.86\%	&	yes	&	0.02\%	\\
p07-s40-v2t2	&	yes	&	2.85\%	&	yes	&	1.01\%	&	yes	&	0.00\%	\\
p07-t40-v2t2	&	No	&	-	&	yes	&	2.49\%	&	yes	&	0.28\%	\\\hline
\end{tabular}
\end{table}  





\begin{comment}

RESULTS FROM PROPOSAL April 2006
We terminated
the branch-and-price algorithm when we found the optimal integer solution or we 
exceeded 8 hours of CPU time.  
%Table \label{BPAlg25} shows the results for the instances with 
%25 customers and table  \label{BPAlg30} shows the results for the instances with 30 customers.
In table  \ref{BPAlg25} and  table \ref{BPAlg30}, we present the results for the instances
with 25 customers and the results for the instances with 30 customers, respectively. The tables include
the LP relaxation bound, the optimal or best integer solution, the integrality gap when the branch-and-price
algorithm terminates, the total number of evaluated nodes in the tree, the total number of columns,
the time spent to solve the pricing problems, the time spent to solve the master problems
and the total time. In addition, the tables report  the number of facilities selected, the number of 
vehicles (pairings) in the integer solution and the average number of customers per vehicle.  

{\bf Instances with 25 customers} We found optimal integer solution for 15 instances (out of 24). The
integrality gaps for the other 6 instances are less than 0.5\%. We have a 30\% gap in one of the
instances and less than 5\% gap in the remaining two instances. 
Except the $p03$ instances, the largest computation time is used for the instances with $v2t2$ 
(the time limit in $t2$ instances is greater than the time limit in $t1$ instances and the vehicle 
capacity in $v2$ instances is larger than the vehicle capacity in $v1$ instances). All of the $p03-f25$ 
instances result in integrality gaps, and among $p03-l25$ instances, $v1t2$ case uses the most  time. 
In general, the problem becomes more difficult as the time limit and the vehicle capacity increases. 

{\bf Instances with 30 customers} We could not find the LP relaxations of the 4 instances within
our time limit. We found optimal or best integer solutions with a gap less than 1\% for 12 instances. 
Our largest integrality gap is 8.2\% for the other instances. All 4 of the instances that could not
be solved at the root node have time limit $t2$. Three of these 4 instances have vehicle capacity $v2$. 
Interestingly, for the $p11$ instances, the  $v1t1$ cases are more difficult to solve than the other $p11$ cases. 
This result shows that the larger vehicle capacity and the time limit are better for $p11$ instances.  

%We have heuristically solved the LRS problem with formulation SPP2. We applied facility 
%location variable branching and total number of vehicles branching. When there is not any
%fractional T and nV variable in the solutions of the leaves of the branch-and-price tree, the
%algorithm terminates and gives the best upper bound as the IP solution. In Table 5, 
%Gap1 represents the \% gap between the LP relaxation and Best IP, Gap2 represents the \% gap 
%between the algorithm's Best IP and optimal IP solution. In most of the instances, the best-IP
%found during the algorithm is equal to the optimal solution of the LRS problem.

\begin{sidewaystable}
%\begin{table}
\centering
\caption{Results of the Branch-and-Price Algorithm for 25 Customers} \label{BPAlg25}
\begin {tabular}[b]{|c|c|c|c|c|c|c|c|c|c|c|c|}
\hline
{\bf Data File} & {\bf LP} & {\bf Optimal or} & {\bf \% Int.} & {\bf \# of BB} &{\bf \# of } &  \multicolumn {3}{|c|}{{\bf Times}} & {\bf \# of } & {\bf \# of } & {\bf Avg. \# of } \\\cline {7-9}
\mbox{}& {\bf Relax.} & {\bf Best IP} & {\bf  Gap} & {\bf Nodes} & {\bf Col.} & {\bf CG} & {\bf LP} & {\bf Total} & {\bf Fac.} & {\bf Veh.}& {\bf Cust./Veh.} \\\hline
p01-f25-v1t1	&	4404.3	&	4618.83	&	0	&	77	&	10992	&	5037.54	&	72.05	&	13407.03	&	2	&	4	&	6.25	\\\hline
p01-f25-v1t2	&	4321.41	&	4406.33* &	0.06	&	31	&	11669	&	21429.9	&	45.56	&	28809.06	&	2	&	3	&	8.33	\\\hline
p01-f25-v2t1	&	4266.38	&	4351.33	&	0	&	53	&	16431	&	13884.42	&	100.67	&	15878.83	&	2	&	3	&	8.33	\\\hline
p01-f25-v2t2	&	4192.23	&	4353*	&	3.52	&	3	&	14037	&	28407.81	&	22.69	&	28962.15	&	2	&	3	&	8.33	\\\hline \hline
p03-f25-v1t1	&	4699.51	&	6702.33* &	29.28	&	253	&	13525	&	1853.45	&	175.16	&	28909.78	&	4	&	6	&	4.17	\\\hline
p03-f25-v1t2	&	4582.2	&	4741*	&	0.18	&	204	&	13308	&	6917.51	&	110.89	&	29171.98	&	2	&	4	&	6.25	\\\hline
p03-f25-v2t1	&	4492.06	&	4674.33	* &	0.17	&	201	&	16559	&	3556.32	&	159.15	&	28901.44	&	2	&	4	&	6.25	\\\hline
p03-f25-v2t2	&	4404.59	&	4676.83 *&	4.07	&	15	&	15243	&	21519.04	&	53.74	&	28826.39	&	2	&	4	&	6.25	\\\hline \hline
p03-l25-v1t1	&	4682.91	&	4721	&	0	&	9	&	9574	&	357.93	&	10.9	&	618.55	&	2	&	4	&	6.25	\\\hline
p03-l25-v1t2	&	4562.92	&	4707.17* &	0.03	&	203	&	13862	&	6108.67	&	89.66	&	28831.01	&	2	&	4	&	6.25	\\\hline
p03-l25-v2t1	&	4526.2	&	4653.83	&	0	&	113	&	12189	&	2416.9	&	65.06	&	6205.48	&	2	&	4	&	6.25	\\\hline
p03-l25-v2t2	&	4416.08	&	4459.33	&	0	&	47	&	15467	&	11040.15	&	59.04	&	18920.63	&	2	&	3	&	8.33	\\\hline \hline
p07-f25-v1t1	&	4560.82	&	4678.5	&	0	&	99	&	13161	&	1775.92	&	74.73	&	5201.75	&	2	&	4	&	6.25	\\\hline
p07-f25-v1t2	&	4457.51	&	4527.5	&	0	&	43	&	13761	&	11206.98	&	49.73	&	18776.02	&	2	&	4	&	6.25	\\\hline
p07-f25-v2t1	&	4408.23	&	4411.83	&	0	&	3	&	12975	&	881.76	&	10.66	&	5585.84	&	2	&	3	&	8.33	\\\hline
p07-f25-v2t2	&	4303.73	&	4411.83* &	0.43	&	23	&	15781	&	28081.55	&	62.56	&	28885.06	&	2	&	3	&	8.33	\\\hline \hline
p11-f25-v1t1	&	7601.51	&	7784	&	0	&	43	&	2055	&	5.68	&	1.51	&	24.01	&	3	&	6	&	4.17	\\\hline
p11-f25-v1t2	&	7521.76	&	7616.5	&	0	&	27	&	2735	&	7.62	&	1.43	&	20.98	&	3	&	6	&	4.17	\\\hline
p11-f25-v2t1	&	7371.89	&	7500.67	&	0	&	33	&	2897	&	9.34	&	2.23	&	24.6	&	3	&	6	&	4.17	\\\hline
p11-f25-v2t2	&	7254.45	&	7459.83	&	0	&	341	&	4786	&	67.75	&	36.03	&	2584.29	&	3	&	5	&	5	\\\hline \hline
p11-l25-v1t1	&	8178.02	&	8443.83	&	0	&	1405	&	1796	&	17.43	&	20.93	&	5100.98	&	4	&	11	&	2.27	\\\hline
p11-l25-v1t2	&	8036.66	&	8231.67	&	0	&	297	&	2038	&	12.47	&	4.39	&	349.94	&	3	&	7	&	3.57	\\\hline
p11-l25-v2t1	&	7633.51	&	7916.67	&	0	&	147	&	1993	&	10.47	&	4.31	&	204.51	&	3	&	6	&	4.17	\\\hline
p11-l25-v2t2	&	7579.66	&	7750.5*	&	0.08	&	2662	&	3752	&	62.27	&	46.43	&	28807.98	&	3	&	6	&	4.17	\\\hline \hline
\multicolumn{12}{l}{* The best IP}
\end{tabular}
%\end{table}
\end{sidewaystable} 



\begin{sidewaystable}
%\begin{table}
\centering
\caption{Results of the Branch-and-Price Algorithm for 30 Customers} \label{BPAlg30}
\begin {tabular}[b]{|c|c|c|c|c|c|c|c|c|c|c|c|}
\hline
{\bf Data File} & {\bf LP} & {\bf Optimal or} & {\bf \% Int.} & {\bf \# of BB} &{\bf \# of } &  \multicolumn {3}{|c|}{{\bf Times}} & {\bf \# of } & {\bf \# of } & {\bf Avg. \# of } \\\cline {7-9}
\mbox{}& {\bf Relax.} & {\bf Best IP} & {\bf  Gap} & {\bf Nodes} & {\bf Col.} & {\bf CG} & {\bf LP} & {\bf Total} & {\bf Fac.} & {\bf Veh.}& {\bf Cust./Veh.} \\\hline
p01-f30-v1t1	&	4821.07	&	4913*	&	0.87	&	101	&	20550	&	16375.4	&	259.83	&	28831.54	&	2	&	7.5	&	6.25	\\\hline
p01-f30-v1t2	&	-	&		&		&	1	&	15840	&	28834.24	&	27.17	&	28873.28	&		&		&		\\\hline
p01-f30-v2t1	&	4582.27	&	4944.67* &	7.26	&	3	&	22584	&	21654.47	&	54.48	&	28948.5	&	2	&	6	&	5	\\\hline
p01-f30-v2t2	&	-	&		&		&	1	&	22160	&	32603.81	&	24.52	&	32636.89	&		&		&		\\\hline \hline
p03-f30-v1t1	&	6467.35	&	6570.17* &	1.06	&	64	&	16568	&	10499.2	&	76.25	&	29837.43	&	4	&	6	&	5	\\\hline
p03-f30-v1t2	&	6370.12	&	6574.33* &	3.11	&	2	&	15165	&	25940.42	&	12.38	&	33184.98	&	3	&	7.5	&	6.25	\\\hline
p03-f30-v2t1	&	6263.17	&	6427.33* &	2.29	&	4	&	22846	&	27591.15	&	37.77	&	29324.21	&	3	&	7.5	&	6.25	\\\hline
p03-f30-v2t2	&	-	&		&		&	1	&	21642	&	30788.36	&	2.62	&	30792.35	&		&		&		\\\hline \hline
p03-l30-v1t1	&	4896.04	&	5014.67* &	0.07	&	95	&	17985	&	5829.87	&	132.57	&	28874.86	&	2	&	6	&	5	\\\hline
p03-l30-v1t2	&	4756.95	&	4905.67* &	2.42	&	33	&	21199	&	25728.99	&	122.48	&	28801.91	&	4	&	4.29	&	3.57	\\\hline
p03-l30-v2t1	&	4679.23	&	4730.17	&	0	&	65	&	21009	&	14066.54	&	154.32	&	28557.37	&	2	&	7.5	&	6.25	\\\hline
p03-l30-v2t2	&	4562.29	&	4965.5*	&	8.12	&	2	&	22270	&	23552.23	&	35.4	&	30840.03	&	2	&	6	&	5	\\\hline \hline
p07-f30-v1t1	&	4765.63	&	4808.33* &	0.05	&	138	&	24113	&	11038.05	&	300.07	&	28859.69	&	2	&	7.5	&	6.25	\\\hline
p07-f30-v1t2	&	4654.19	&	4882.33* &	4.55	&	6	&	20532	&	21641.46	&	49.52	&	28928.4	&	2	&	7.5	&	6.25	\\\hline
p07-f30-v2t1	&	4558.21	&	4729.17* &	2.47	&	17	&	23982	&	25993.7	&	129.95	&	28870.75	&	2	&	7.5	&	6.25	\\\hline
p07-f30-v2t2	&	-	&		&		&	1	&		&	29020.8	&	31.9	&	29059.58	&		&		&		\\\hline \hline
p11-f30-v1t1	&	7916.08	&	8058*	&	0.29	&	1960	&	4089	&	68.31	&	64.82	&	28816.77	&	3	&	4.29	&	3.57	\\\hline
p11-f30-v1t2	&	7791.1	&	8036.33* &	0.34	&	1430	&	4893	&	73.86	&	54.74	&	28843.89	&	3	&	4.29	&	3.57	\\\hline
p11-f30-v2t1	&	7531.81	&	7634.67	&	0	&	137	&	3670	&	32.27	&	6.56	&	585.2	&	3	&	5	&	4.17	\\\hline
p11-f30-v2t2	&	7469.84	&	7634.67	&	0	&	309	&	5376	&	86.97	&	18.05	&	2743.63	&	3	&	5	&	4.17	\\\hline \hline
p11-l30-v1t1	&	9929.75	&	10043	&	0	&	525	&	2729	&	24.61	&	7.18	&	2973.63	&	4	&	4.29	&	3.57	\\\hline
p11-l30-v1t2	&	9837.09	&	9937.33	&	0	&	49	&	3182	&	18.53	&	2.13	&	73.18	&	4	&	4.29	&	3.57	\\\hline
p11-l30-v2t1	&	9645.66	&	9758.17	&	0	&	127	&	3345	&	34.53	&	6.82	&	331.64	&	4	&	4.29	&	3.57	\\\hline
p11-l30-v2t2	&	9416.78	&	9500.67	&	0	&	105	&	4488	&	55.82	&	5.75	&	324.02	&	4	&	5	&	4.17	\\\hline \hline
\multicolumn{12}{l}{- The algorithm cannot find a solution within the time limit}\\
\multicolumn{12}{l}{* The best IP}
\end{tabular}
%\end{table}
\end{sidewaystable} 
\end{comment}
